\documentclass[11pt]{article}
\usepackage[a4paper,margin=25mm]{geometry}
\usepackage{amsmath,amssymb,amsthm,booktabs,array}
\usepackage[T1]{fontenc}
\usepackage{lmodern}
\usepackage{xurl}
\usepackage[colorlinks=true,urlcolor=blue,citecolor=blue,linkcolor=blue]{hyperref}
\newtheorem{theorem}{Theorem}
\newtheorem{lemma}{Lemma}
\newtheorem{corollary}{Corollary}
\newcommand{\N}{\mathbb{N}_0}
\newcommand{\lpf}{\operatorname{lpf}}
\title{Exact least prime factors and exceptional classes for Wu's conjectures on\protect\\
$(2^{2^n}+1)2^{2^{2^n}}-1$}
\author{Haobo Ma\\\small ChronoAI Pte Ltd, Singapore
\and Wenlin Zhang\\\small National University of Singapore\\\small ChronoAI Pte Ltd, Singapore}
\date{13 November 2026}
\begin{document}
\maketitle
\begin{abstract}
For Chai Wah Wu's Conjectures 1--4 in OEIS A227244, revision 54, FORMULA
(10 July 2026), we extend the bounds proved by Schott (23 October 2019)
and Wu (10 July 2026)~\cite{oeis} to exact least-prime-factor formulas,
including all excluded classes.
The values are 101 on $20n+16$, 47 on $220n+8$, and 149 on $72n+32$
outside $n\equiv2\pmod5$ and $n\equiv18\pmod{55}$, where they are 101
and 47, respectively. On $168n+43$ the value is 59 except at
$n\equiv26\pmod{55}$, where it is 47.
A periodicity lemma and finite modular tables prove these formulas.
We also characterize all indices with value 11 and extend Fried's (2025)
Fermat-number counterexample to an exact criterion for Cloitre's totient
condition~\cite{fried}.
\end{abstract}

\section{Statements and the noncoverage condition}
For $m\in\N$, put
\[
 G_m=(2^{2^m}+1)\,2^{2^{2^m}}-1,
 \qquad a(m)=\lpf(G_m),
\]
where $\lpf$ denotes the least prime divisor of an integer greater than one.
Since $G_0=11$ and $G_m\geq 11$, this is defined for every $m$.
Wu's four conjectures concern the values of this sequence along arithmetic
progressions~\cite{oeis}. The issue is minimality, rather than divisibility
alone: the same integer can have several small prime divisors.
Schott's recorded proof excludes the primes below 11; Wu's proofs give
$a(4n+5)=11$ and $a(20n+16)\leq101$. The equalities below add the
remaining smaller-prime exclusions and specify the omitted classes.

\begin{theorem}\label{thm:main}
For every $n\in\N$,
\begin{align}
 a(20n+16)&=101,\label{eq:c1}\\
 a(220n+8)&=47,\label{eq:c2}\\
 a(72n+32)&=
 \begin{cases}
 101,&n\equiv2\pmod5,\\
 47,&n\equiv18\pmod{55},\\
 149,&n\not\equiv2\pmod5\text{ and }n\not\equiv18\pmod{55},
 \end{cases}\label{eq:c3}\\
 a(168n+43)&=
 \begin{cases}
 47,&n\equiv26\pmod{55},\\
 59,&n\not\equiv26\pmod{55}.
 \end{cases}\label{eq:c4}
\end{align}
The two exceptional cases in~\eqref{eq:c3} are disjoint.
\end{theorem}

We read Conjecture 3's noncoverage clause as exclusion of the index from
the first two progressions. This condition is essential. Indeed,
\begin{align*}
 72n+32\in\{20k+16:k\in\N\}&\iff n\equiv2\pmod5,\\
 72n+32\in\{220k+8:k\in\N\}&\iff n\equiv18\pmod{55}.
\end{align*}
To check the forward implications, reduce the two equations modulo 20 and 220,
respectively: $12n\equiv4\pmod{20}$ and $72n\equiv196\pmod{220}$.
For the converses, $n=5h+2$ gives $k=18h+8$, while $n=55h+18$ gives $k=18h+6$.
These are nonnegative integers. The two conditions cannot coincide because
$18\equiv3\pmod5$.
The unconditional equality, which is not asserted in the source, fails:
$n=2$ gives $a(176)=101$, and $n=18$ gives $a(1328)=47$.
Theorem~\ref{thm:main} proves the conjectured assertion on the specified
remaining indices and determines both excluded sets.

\section{A finite reduction for the nested powers}
The following elementary lemma supplies the all-index reduction; modular
arithmetic never requires constructing $G_m$.

\begin{lemma}\label{lem:reduction}
Let $q$ be an odd prime and $d=q-1$. Let $s,t,p,L$ be positive integers
satisfying
\begin{align}
 2^d&\equiv1\pmod q,&
 2^{2p}&\equiv2^p\pmod d,\label{eq:period1}\\
 2^{sL+t}&\equiv2^t\pmod d,&
 2^{sL+t}&\equiv2^t\pmod p,&p&\leq2^t.\label{eq:period2}
\end{align}
For $0\leq r<L$, define
\begin{equation}\label{eq:residue}
 R_{q,s,t}(r)=
 \left(2^{\,2^{sr+t}\bmod d}+1\right)
 2^{\,2^{\,p+(2^{sr+t}\bmod p)}\bmod d}-1\pmod q.
\end{equation}
Then $G_{sn+t}\equiv R_{q,s,t}(n\bmod L)\pmod q$ for all $n\in\N$.
\end{lemma}

\begin{proof}
If $2^{sL+t}\equiv2^t\pmod M$, multiplication by $2^{sk}$ shows that
$2^{s(k+L)+t}\equiv2^{sk+t}\pmod M$ for every $k\geq0$.
Consequently $2^{sn+t}\equiv2^{sr+t}$ modulo both $d$ and $p$, where
$r=n\bmod L$.
Similarly, the second congruence in~\eqref{eq:period1}, multiplied by $2^k$,
shows that $2^{k+2p}\equiv2^{k+p}\pmod d$ for all $k\geq0$.
Thus, for $e\geq p$,
\[
 2^e\equiv2^{p+(e\bmod p)}\pmod d.
\]
Take $m=sn+t$ and $e=2^m\geq2^t\geq p$. It follows that
\[
 2^{2^m}\equiv2^{p+(2^{sr+t}\bmod p)}\pmod d.
\]
Finally $2^d\equiv1\pmod q$ permits reduction modulo $d$ of both exponents
in $(2^{2^m}+1)2^{2^{2^m}}-1$.
The first exponent reduces to $2^{sr+t}\bmod d$; the second reduces to
$2^{p+(2^{sr+t}\bmod p)}\bmod d$. This is~\eqref{eq:residue}.
\end{proof}

\subsection*{The modular tables}
Table~\ref{tab:parameters} gives $p$ and $L_j$ for each required odd prime,
with $d=q-1$ and
\[
 (s_1,t_1)=(20,16),\quad(s_2,t_2)=(220,8),\quad
 (s_3,t_3)=(72,32),\quad(s_4,t_4)=(168,43).
\]
For every populated column, direct modular exponentiation verifies all the
premises of Lemma~\ref{lem:reduction}. The zero sets of~\eqref{eq:residue}
are exactly those in Table~\ref{tab:zeros}; every other populated row has
empty zero set. Primes above the indicated column's target are not needed.
The supplied Python program \texttt{verify.py} checks all parameters and zero
sets in \texttt{certificates.json}. These finite checks enumerate residue
classes rather than sample indices.

\begin{table}[ht]
\centering
\begin{tabular}{cccl}
\toprule
Progression & $q$ & $L$ & $\{r:0\leq r<L,\ R_{q,s,t}(r)=0\}$\\
\midrule
$20n+16$ &101&1&$\{0\}$\\
$220n+8$ &47&1&$\{0\}$\\
$72n+32$ &47&55&$\{18\}$\\
           &101&5&$\{2\}$\\
           &149&1&$\{0\}$\\
$168n+43$ &47&55&$\{26\}$\\
            &59&1&$\{0\}$\\
\bottomrule
\end{tabular}
\caption{All nonempty zero sets in the modular tables.}\label{tab:zeros}
\end{table}

\begin{proof}[Proof of Theorem~\ref{thm:main}]
Every $G_m$ is odd, so 2 is never a prime divisor.
For the first progression, the table excludes every odd prime less than 101
and proves divisibility by 101. Thus its least prime divisor is 101.
The identical argument with target 47 proves~\eqref{eq:c2}.

In the third progression, 149 always divides $G_{72n+32}$.
Among primes below 149, only 47 and 101 can divide it.
The former divides precisely when $n\equiv18\pmod{55}$, and the latter
precisely when $n\equiv2\pmod5$. Their classes are disjoint, so the least
prime divisor is 47, 101 or 149 exactly as in~\eqref{eq:c3}.
For the fourth progression, 59 always divides $G_{168n+43}$, and the only
possible smaller prime divisor is 47, occurring precisely when
$n\equiv26\pmod{55}$. This proves~\eqref{eq:c4}, including its exceptional case.
\end{proof}

\section{Short consequences and the Fermat-number cross-reference}
The earlier A227244 formulas and its cross-reference A000215 have the
following elementary consequences. The first lower bound and the forward
implication in~\eqref{eq:eleven} reprove Schott's and Wu's recorded
results~\cite{oeis}; the reverse implication completes their index description.

\begin{corollary}\label{cor:eleven}
For every $m\in\N$, $a(m)\geq11$, and
\begin{equation}\label{eq:eleven}
 a(m)=11\quad\Longleftrightarrow\quad
 m=0\ \text{or}\ (m\geq5\text{ and }m\equiv1\pmod4).
\end{equation}
\end{corollary}
\begin{proof}
The integers $G_m$ are odd. For $m\geq1$, both powers of 2 in their
definition are 1 modulo 3, so $G_m\equiv1\pmod3$.
For $m\geq2$, both are also 1 modulo 5, so $G_m\equiv1\pmod5$.
Modulo 7 the first power is 2 or 4, hence its successor is 3 or 5;
the second power is in $\{1,2,4\}$. None of the resulting products is
1 modulo 7. The remaining cases are $G_0=11$ and $G_1=79$.
This excludes every prime below 11.

For $m\geq2$, the residues of $2^m$ modulo 10, in the order
$m\equiv0,1,2,3\pmod4$, are $6,2,4,8$. Also
$2^{2^m}\equiv6\pmod{10}$, so
$2^{2^{2^m}}\equiv2^6\equiv9\pmod{11}$.
Consequently the residues of $G_m$ modulo 11 are $1,0,9,2$ in that order.
Together with $G_0=11$ and $G_1=79$, this proves~\eqref{eq:eleven}.
\end{proof}

Put $F_j=2^{2^j}+1$, as in A000215~\cite{fermat}.
Schott's identity $G_m=F_m2^{F_m-1}-1$ is recorded in A227244, COMMENTS
(23 October 2019). Its relation to earlier Fermat numbers is also exact.
\begin{corollary}\label{cor:coprime}
For $0\leq j<m$, $G_m\equiv1\pmod{F_j}$, whereas
$G_m\equiv-1\pmod{F_m}$. In particular,
$\gcd(G_m,\prod_{j=0}^mF_j)=1$.
\end{corollary}
\begin{proof}
Since $2^{2^j}\equiv-1\pmod{F_j}$, powers of 2 modulo $F_j$ have
period dividing $2^{j+1}$. For $j<m$, both exponents $2^m$ and
$2^{2^m}$ are divisible by $2^{j+1}$. Thus the two powers in $G_m$
are 1 modulo $F_j$. The congruence modulo $F_m$ follows from its
factorization $G_m=F_m2^{F_m-1}-1$.
\end{proof}

Fried~\cite{fried} proves Layman's binary-to-quaternary formula and the
Fermat-number direction of Hasler's conjecture, and gives $F_5$ as a
counterexample to Cloitre's proposed primality criterion. The next corollary
reproves the first two statements and strengthens the description of the third.
Here $T(\sum b_i2^i)=\sum b_i4^i$ for binary digits $b_i$, and $\varphi$
is Euler's totient function, with $\varphi(1)=1$.
\begin{corollary}\label{cor:fermat}
For $j\in\N$, $T(F_j)=F_{j+1}$. For $j\geq2$,
$(F_j-1)^4-1$ divides $2^{F_j-1}-1$. Moreover,
\begin{equation}\label{eq:totient}
 \varphi(F_j-2)=(F_j-1)/2\quad\Longleftrightarrow\quad 0\leq j\leq5.
\end{equation}
In general this totient equality holds exactly when all $F_i$ with $i<j$
are prime. The implication from primality to this equality for every $j$
is therefore equivalent to the assertion that no Fermat prime has index
$j\geq6$.
\end{corollary}
\begin{proof}
The binary nonzero positions of $F_j$ are $2^j$ and 0, so
$T(F_j)=4^{2^j}+1=F_{j+1}$.
For $j\geq2$, $j+2\leq2^j$, and hence
$2^{j+2}\mid2^{2^j}$. The geometric-sum identity
$2^u-1\mid2^v-1$ when $u\mid v$ gives the stated divisibility.

Successive differences of squares give
$F_j-2=\prod_{i<j}F_i$ (including the empty product for $j=0$).
Also $F_k\equiv2\pmod{F_i}$ for $i<k$, so these odd numbers are
pairwise coprime. Multiplicativity of $\varphi$ therefore yields
\[
 \varphi(F_j-2)=\prod_{i<j}\varphi(F_i)
 \leq\prod_{i<j}(F_i-1)=2^{2^j-1}=(F_j-1)/2.
\]
Each factor is positive, and $\varphi(N)=N-1$ for $N>1$ exactly when
$N$ is prime. Equality holds precisely when every earlier $F_i$ is prime.
The numbers $3,5,17,257,65537$ are prime by trial division, whereas
$F_5=641\cdot6700417$ is composite. This proves~\eqref{eq:totient}
and the final equivalence.
\end{proof}

\section{Coverage and verification scope}
The source statements have the following coverage; extra consequences are
identified separately from what the sources assert.
\begin{center}
\begin{tabular}{>{\raggedright\arraybackslash}p{35mm}>{\raggedright\arraybackslash}p{104mm}}
\toprule
Source item & Mathematical coverage\\
\midrule
Conjectures 1 and 2 & Proved by~\eqref{eq:c1}--\eqref{eq:c2}.\\
Conjecture 3 & Proved with progression noncoverage; its excluded classes
are determined in~\eqref{eq:c3}.\\
Conjecture 4 & Proved and extended to its excluded class in~\eqref{eq:c4}.\\
Earlier FORMULA block & Reproved; prior proofs Schott (2019), Wu (2026).
Corollary~\ref{cor:eleven} also gives the converse, not posed in the source.\\
COMMENTS, upper bound & Prior proof Wu (2026); equation~\eqref{eq:c1}
strengthens the bound to equality.\\
COMMENTS, individual terms & $a(16)=a(36)=101$, $a(32)=149$ and $a(43)=59$
follow from Theorem~\ref{thm:main}; the other large terms and historical
primality claims are not covered.\\
A000215, Layman & Reproved; prior proof Fried (2025), first paragraph.\\
A000215, Hasler & Fermat-number implication reproved; prior proof Fried
(2025), Theorem 1. Converse not covered.\\
A000215, Cloitre & Refuted-and-corrected: Fried (2025), last paragraph,
gives the composite $F_5$; Corollary~\ref{cor:fermat} gives the exact
replacement criterion~\eqref{eq:totient}.\\
A000215, other conjectures & The five-Fermat-prime assertion and
Wesolowski's least-factor bound are not covered.\\
Additional consequences & Corollary~\ref{cor:coprime} and the full totient
criterion are not posed in the sources; proved here as written consequences.\\
\bottomrule
\end{tabular}
\end{center}
The COMMENTS also give Wu's proof of $a(20n+16)\leq101$; equation~\eqref{eq:c1}
adds the matching lower bound.

The written proof consists of the periodicity lemma, the interpretation of
the two noncoverage conditions, and the deduction of least-prime-factor values
from the finite zero sets. The finite computations establish the modular
premises and all the zero sets in Tables~\ref{tab:zeros} and~\ref{tab:parameters}.
The following pinned Lean declarations express the corresponding prime,
divisibility and lower-bound statements at source revision
\texttt{ec1e94b3ab1ae33ea4f94d5484df26d748e2175e}:
\begin{itemize}
\item \href{https://github.com/the-omega-institute/trureturing/blob/ec1e94b3ab1ae33ea4f94d5484df26d748e2175e/D5/S3/Arith/Congruence/WuA227244C1Equality.lean}{\texttt{WuA227244C1Equality.result}}:
101 is prime, divides $G_{20n+16}$, and is no larger than any prime divisor.
\item \href{https://github.com/the-omega-institute/trureturing/blob/ec1e94b3ab1ae33ea4f94d5484df26d748e2175e/D5/S3/Arith/Congruence/WuA227244C2Equality.lean}{\texttt{WuA227244C2Equality.result}}:
the analogous three assertions for 47 on $220n+8$.
\item \href{https://github.com/the-omega-institute/trureturing/blob/ec1e94b3ab1ae33ea4f94d5484df26d748e2175e/D5/S3/Arith/Primes/WuA227244ConjectureThree.lean}{\texttt{WuA227244ConjectureThree.result}}:
the analogous assertions for 149 on $72n+32$, assuming both progression
noncoverage conditions.
\item \href{https://github.com/the-omega-institute/trureturing/blob/ec1e94b3ab1ae33ea4f94d5484df26d748e2175e/D5/S3/Arith/Primes/WuA227244ConjectureFour.lean}{\texttt{WuA227244ConjectureFour.result}}:
the analogous assertions for 59 on $168n+43$, assuming $n\bmod55\ne26$.
\end{itemize}
The exceptional-class equalities in~\eqref{eq:c3}--\eqref{eq:c4} are consequences
of the written argument and finite tables; the four declarations above have
the stated, narrower scopes. Corollaries~\ref{cor:eleven}--\ref{cor:fermat}
are written proofs; the accompanying checker also verifies their finite
residue tables, the five small primes and the factorization of $F_5$.

\begin{thebibliography}{9}
\bibitem{oeis}
The OEIS Foundation Inc., \emph{A227244: Smallest prime factor of
$(2^{2^n}+1)2^{2^{2^n}}-1$}, revision 54, 12 July 2026.
FORMULA, Chai Wah Wu, 10 July 2026; COMMENTS, Bernard Schott, 23 October 2019,
and Chai Wah Wu, 10 July 2026.
\url{https://oeis.org/A227244}.
\bibitem{fermat}
The OEIS Foundation Inc., \emph{A000215: Fermat numbers},
revision 344, 23 August 2026.
\url{https://oeis.org/A000215}. Accessed 13 November 2026.
\bibitem{fried}
Sela Fried, \emph{On several conjectures stated in A000215}, 2025,
pp.~1--2; linked from A000215, COMMENTS (31 October 2025).
\url{https://oeis.org/A000215/a000215_1.pdf}.
\end{thebibliography}
\noindent\small Copyright \copyright\ 2026 Haobo Ma and Wenlin Zhang.
Licensed under Creative Commons Attribution 4.0 International (CC BY 4.0),
\url{https://creativecommons.org/licenses/by/4.0/}.

\clearpage
\appendix
\section{Parameters for the finite modular proof}
\begin{table}[ht]
\centering
\begin{tabular}{rrrrrr@{\hspace{9mm}}rrrrrr}
\toprule
$q$ & $p$ & $L_1$ & $L_2$ & $L_3$ & $L_4$ & $q$ & $p$ & $L_1$ & $L_2$ & $L_3$ & $L_4$ \\
\midrule
3 & 1 & 1 & 1 & 1 & 1 & 67 & 10 & 1 & -- & 5 & -- \\
5 & 2 & 1 & 1 & 1 & 1 & 71 & 12 & 3 & -- & 1 & -- \\
7 & 2 & 1 & 1 & 1 & 1 & 73 & 6 & 3 & -- & 1 & -- \\
11 & 4 & 1 & 1 & 1 & 1 & 79 & 12 & 3 & -- & 1 & -- \\
13 & 2 & 1 & 1 & 1 & 1 & 83 & 20 & 1 & -- & 5 & -- \\
17 & 4 & 1 & 1 & 1 & 1 & 89 & 10 & 1 & -- & 5 & -- \\
19 & 6 & 3 & 3 & 1 & 1 & 97 & 6 & 1 & -- & 1 & -- \\
23 & 10 & 1 & 1 & 5 & 5 & 101 & 20 & 1 & -- & 5 & -- \\
29 & 3 & 3 & 3 & 1 & 1 & 103 & 8 & -- & -- & 1 & -- \\
31 & 4 & 1 & 1 & 1 & 1 & 107 & 52 & -- & -- & 13 & -- \\
37 & 6 & 3 & 3 & 1 & 1 & 109 & 18 & -- & -- & 1 & -- \\
41 & 4 & 1 & 1 & 1 & 1 & 113 & 6 & -- & -- & 1 & -- \\
43 & 6 & 3 & 3 & 1 & 1 & 127 & 6 & -- & -- & 1 & -- \\
47 & 11 & 11 & 1 & 55 & 55 & 131 & 12 & -- & -- & 1 & -- \\
53 & 12 & 3 & -- & 1 & 1 & 137 & 8 & -- & -- & 1 & -- \\
59 & 28 & 21 & -- & 7 & 1 & 139 & 22 & -- & -- & 55 & -- \\
61 & 4 & 1 & -- & 1 & -- & 149 & 36 & -- & -- & 1 & -- \\
\bottomrule
\end{tabular}

\caption{Parameters in Lemma~\ref{lem:reduction}, with $d=q-1$.
Column $j$ uses $(s_j,t_j)$ specified in the text. A dash denotes a prime
above the target for that column. The zero sets are in Table~\ref{tab:zeros}.}
\label{tab:parameters}
\end{table}
\end{document}
